\textbf{Decision Gate Metrics and Go/No-Go Criteria}

Phase~I success will be evaluated through quantitative metrics
demonstrating predictive accuracy, computational efficiency,
and measurable AI advantage relative to conventional
high-fidelity CFD workflows using Nek5000/RS.

\textbf{1. Integrated Data and Latent Representation.}
A unified geometry-aware dataset combining simulation
and experimental thermal–fluid data will be constructed.
Success criteria include:

\begin{itemize}[noitemsep]

\item Successful integration of $\geq 95\%$ of selected datasets
into a consistent latent representation.

\item Reconstruction fidelity of encoded flow fields with
relative $L_2$ error $\leq 3\%$ for representative
thermal–fluid variables.

\end{itemize}

\textbf{2. Predictive Accuracy and Generalization.}
The foundation model will be validated across unseen
geometries and operating regimes. Success requires:

\begin{itemize}[noitemsep]

\item Prediction error $\leq 5\%$ for engineering quantities
including pressure drop and heat-transfer coefficients.

\item Robust performance across at least three geometry
classes and multiple Reynolds-number regimes.

\end{itemize}

\textbf{3. Design Acceleration and Optimization Efficiency.}
AI-enabled latent-space optimization will be compared
to traditional brute-force parameter sweeps. Success
criteria include:


\begin{itemize}[noitemsep]

\item Reduction of required high-fidelity CFD evaluations
by $\geq 5$--$10\times$.

\item Demonstration of stable multi-objective optimization
balancing thermal performance and pressure-drop penalties.

\end{itemize}

\textbf{4. Demonstration of AI Advantage Through Scaling.}
AI advantage will be quantified using scaling-based
performance metrics:

\begin{itemize}[noitemsep]

\item Reduction in prediction error of $\geq 30\%$
when training dataset size is increased by $2\times$.

\item Efficient GPU scaling demonstrating proportional
reduction in training or inference time with increasing
computational resources.

\item Demonstration of surrogate inference throughput
$\geq 100\times$ faster than equivalent CFD evaluation.

\end{itemize}

\textbf{Go/No-Go Criterion.}
Advancement to Phase~II requires simultaneous
demonstration of:

\begin{itemize}[noitemsep]

\item Prediction accuracy $\leq 5\%$ across unseen geometries,

\item $\geq 5\times$ reduction in required CFD evaluations,

\item Measurable scaling-based AI advantage,

\item $\geq 10\times$ reduction in end-to-end
design-cycle time relative to traditional workflows.

\end{itemize}

Meeting these criteria will validate the feasibility of
foundation-model–driven workflows for accelerated
thermal–fluid design and justify Phase~II expansion.
